% Generated draft: requires mathematical and source review.
\documentclass[11pt]{article}
\usepackage[margin=25mm]{geometry}
\usepackage{fontspec}
\setmainfont{Latin Modern Roman}
\usepackage{amsmath,amssymb,mathtools,unicode-math}
\setmathfont{Latin Modern Math}
\usepackage{enumitem,needspace,xurl,hyperref}
\setlength{\parindent}{0pt}
\setlength{\parskip}{5pt}
\setlength{\emergencystretch}{3em}
\newcommand{\N}{\mathbb{N}}
\newcommand{\Z}{\mathbb{Z}}
\newcommand{\Q}{\mathbb{Q}}
\newcommand{\R}{\mathbb{R}}
\newcommand{\C}{\mathbb{C}}
\newcommand{\F}{\mathbb{F}}
\newcommand{\A}{\mathbb{A}}
\newcommand{\PP}{\mathbb{P}}
\newcommand{\E}{\mathbb{E}}
\newcommand{\eps}{\varepsilon}
\newcommand{\dd}{\,\mathrm{d}}
\begin{document}
\section*{UnsolvedMath ID: 30003418}
\textbf{Source problem code:} catalogue entry
\subsection{Definitions and mathematical statement}
The supplied record asks about the assertion
\[
\mathsf{A}:\qquad E_{\omega_1}\leq_{a\Delta^1_1}F_{\omega_1}.
\]
Here \(\omega_1\) conventionally denotes the first uncountable ordinal. The record does not define the equivalence relations \(E_{\omega_1}\) and \(F_{\omega_1}\), their underlying spaces, or the relation \(\leq_{a\Delta^1_1}\). Consequently, \(\mathsf{A}\) is a named assertion from the source, not a fully specified assertion in this draft.

For context, if \(E\) and \(F\) are equivalence relations on sets \(X\) and \(Y\), respectively, an ordinary reduction is a function \(f:X\to Y\) satisfying
\[
\forall u,v\in X\quad \bigl(u\,E\,v\iff f(u)\,F\,f(v)\bigr).
\]
This definition does not specify the source's additional requirements: the meaning of ``almost,'' the effective definability requirement, and the permitted real parameters must be supplied before it can define \(\leq_{a\Delta^1_1}\).

\textbf{Question as supplied.} Determine the consistency strength of \(\mathsf{A}\). Does \(\mathsf{A}\) imply the existence of a real whose admissibility spectrum is sufficiently sparse that its admissible ordinals are cardinals of \(L\), and from which \(0^\#\) can be recovered?

Here a real may conventionally be represented by a subset of \(\omega\), and \(L\) denotes the constructible universe. For an ordinal \(\alpha\), saying that \(\alpha\) is an \(L\)-cardinal means
\[
L\models\text{``}\alpha\text{ is a cardinal.''}
\]
The notation \(0^\#\) denotes zero sharp, when it exists. The record does not specify the definition or ordinal range of the admissibility spectrum, a quantified sparsity condition, or the meaning of recovery of \(0^\#\). These phrases therefore remain unresolved parts of the question. In particular, the recovery requirement cannot silently be replaced by the assertion that \(0^\#\) exists. The background theory and intended comparison of consistency strengths are also unspecified. These omissions prevent a complete formal statement without additional source evidence.
\subsection{Short English statement}
The problem asks for the consistency strength of an almost-\(\Delta^1_1\) reduction from \(E_{\omega_1}\) to \(F_{\omega_1}\), and whether such a reduction forces a real with a sparse admissibility spectrum consisting of \(L\)-cardinals and permitting recovery of \(0^\#\). The supplied record leaves the reduction notion and several essential qualifications undefined, so this is a statement-only draft requiring human mathematical review.
\subsection{Sources}
\begin{itemize}
\item Local UnsolvedMath registry, permanent ID \(30003418\), title: \emph{Consistency Strength of Almost Effective Equivalence Reductions}; source problem code \texttt{catalogue entry}. Collection: \emph{Oberwolfach Reports Open Problems}. The record gives proposed year 2017 and names no proposer.
\item Supplied report citation: \emph{Set Theory (2018)}, Oberwolfach Reports, DOI \texttt{10.4171/owr/2017/11}, pp.~527--589. \url{https://doi.org/10.4171/owr/2017/11}. The supplied title year is 2018, while the DOI and proposed-year metadata refer to 2017; this discrepancy is unresolved. No exact problem page is supplied.
\item Supplied literature reference: \emph{Countable admissible ordinal equivalence relation}. \url{https://arxiv.org/abs/1601.07924}.
\item Supplied literature reference: \emph{Author version of the admissibility paper}. \url{https://williamchan-math.github.io/papers/countable_admissible_ordinal_equivalence_relation.pdf}.
\item Literature assessment dated exactly \texttt{2026-08-21}: the supplied curation labels the problem open and reports non-reducibility in \(L\), positive almost reductions from a measurable cardinal, and an undetermined exact consistency strength. This is supplied dated curation, not an independently verified assessment. The background explicitly requests independent verification of time-sensitive status and cited claims.
\item The registry labels the statement \texttt{corrected\_verified}, explaining that notation, typography, and sentence structure were repaired while preserving mathematical content. This supplied editorial label does not establish mathematical correctness or resolve the missing definitions. The references above were supplied with the record and have not been independently consulted for this draft.
\end{itemize}
\end{document}
